\section{Exchange-correlation functionals}
\subsection{Local density approximation (LDA)}
We don't know how to write the exchange correlation term $E_{XC}[\rho]$ using $\rho$, so we need to write it for the electron gas, and assume it works for all systems.
%\begin{align}
%    E_{XC}[\rho]=-\frac{3N}{4}\left(\frac{3\rho}{\pi}\right)^{1/3}
%\end{align}

We do not have shell structure(s,p,d) if we calculate $T^{\sigma}$ using Orbital-free DFT. The term does not consider Coulomb interaction. However, It works relatively well for "warm dense materials" where temperature is larger than $10,000K$, because the effect of orbital is less.

In this approach, the local exchange energy is:
\begin{align}
    E_{xc}[\rho] &= \int d\vec{r} \rho (\vec{r}) \varepsilon_{xc} ([\rho] , \vec{r})\\
    \varepsilon_{xc}([\rho],\vec{r}) &= \varepsilon_x([\rho],\vec{r})+\varepsilon_c([\rho],\vec{r})
\end{align}

The spin component can be written as:
\begin{align}
    E^{LDA}_{XC}[\rho^{\uparrow},\rho^{\downarrow}]=\int d\vec{r}\rho(\vec{r})\varepsilon^V_{XC}([\rho^{\uparrow},\rho^{\downarrow}],\vec{r})
\end{align}

\begin{align}
    V_{XC} = \frac{\delta E_{XC}[\rho]}{\delta\rho} = \varepsilon_{xc}([\rho],\vec{r})+\rho(\vec{r}) \frac{\partial \varepsilon_{xc}([\rho],\vec{r})}{\partial \rho}
\end{align}

\subsection{Generalized-gradient approximation(GGA)}
\begin{align}
    E^{GGA}_{xc}[\rho^{\uparrow},\rho^{\downarrow}] = \int d\vec{r} \rho(\vec{r}) \varepsilon_{xc} (\rho^{\uparrow}, \rho^{\downarrow}, \vert \nabla \rho^{\uparrow}\vert , \vert \nabla \rho ^{\downarrow} \vert , r)
\end{align}

There are several variations of GGA, include PBE and BLYP.

\subsection{Meta GGA}
"SCAN", $$\tau (\vec{r}) := \sum_{i} \frac{1}{2} \vert \nabla \varphi_{i} ( \vec{r}) \vert ^2$$\\
$$E^{meta GGA}_{XC}[\rho] = \int d\vec{r}\rho(\vec{r}) \varepsilon_{xc}(\rho,\vert \nabla \rho \vert , \tau, \vec{r})$$

\subsection{Hybrid functional, orbital dependent functional}
"self-interaction error": in hartree, cannot remove electron self-interaction from the integral

\begin{enumerate}
    \item PBEO: $E_{xc} = \frac{1}{4}E^{HF}_x + \frac{3}{4}E^{PBE}_{x} + E^{PBE}_c$
    \item B3LYP: $E_{xc} = E^{LDA}_x+a_0(E^{HF}_x-E^{LDA}_x)+a_x(E^{Becke}_x-E^{LDA}_x)+E^{LDA}_c+a_c(E^{Becke}-E^{LDA}))$
\end{enumerate}

\subsection{self-consistent hybrid functional for condensed systems}
PRB 89, 195112(2014)

\begin{align}
    E_{xc} &= \alpha E^{HF}_{x} + (1-\alpha) E^{GGA}_x + E^{GGA}_c\\
    \alpha &= \frac{1}{\varepsilon_{\infty}}\\
    \vec{D} &= \varepsilon_{\infty} \vec{E}\\
    V(r) &= k \frac{q_1q_2}{r} \implies V(r) = k \frac{1}{\varepsilon_{\infty}} \frac{q_1 q_2}{r}
\end{align}

Physical meaning of $\varepsilon_{\infty}$: the screening effect by electrons only.

Guess $\alpha \rightarrow$ DFT $\displaystyle\rightarrow \alpha = \frac{1}{\varepsilon_{\infty}} \rightarrow$ guess $\alpha$

\subsection{Range-separated hybrid functionals}
PRB93, 235106(2016)

\begin{align}
    E_{xc} = \alpha E^{HF, lr}_{x} + \beta E^{HF, sr}_{x} + (1-\alpha) E^{GGA, lr} _{x} + (1-\beta)E^{GGA, sr}_{x} + E^{GGA}_c
\end{align}
For a long range, we need to apply a tanh-like error function: $\text{erf}(\mu\vert\vec {r}=\vec{r'}\vert))$

Notice that for PBEO: $\alpha = 0.25$, $\beta = 0$, $\mu \rightarrow \infty$; for HSE06: $\alpha = 0.25$, $\beta = 0$, $\mu = 0.11$

\subsection{van der Waals functional}
\begin{enumerate}
    \item Grimme: DFT-(D2/D3/...) obtain the coefficient $C_6$ for the interaction $-\frac{C_6}{r^6}$ from a table
    \item Tkatchenko-Scheffeler theoretically calculate (to a certain extent) $C_6$ (PRL102, 073005 (2009))
    \item DRSLL: PRL 90 2464401(2004) More about random phase
\end{enumerate}
%wtf josun
 
